% Nouveau EGA V (cote 25): archive cote 25, site page 171 (PDF page 172).
% typescript. Draft compared: none.
% AI-transcribed by Claude Opus 5.5 (claude-opus-5-5, adaptive thinking, effort high), Message Batches API, batch ega25b (msgbatch_01A7n5YuXubUPNK1zsepi8RG), 2026-11-13; not yet checked by eye.
\scanpage{25/171}
\mspage{6}

sa\over{\unsure{i}}{n}t qu'en tout point de $u(V^{o})$, $Z$ et $Z^{o}$ sont lisses sur $M$, $Z^{o}$ étant de codi
mension $m$ dans $Z$ en \over{s}{c}es points. Ceci implique automatiquement que $V^{o}$ est
loc de présentation finie sur $Y$, et que $\underline{N}_{V^{o}/X}$ est l'image de
\xout{\unsure{xxx}}
\[ \underline{E}^{o} = u\,(\underline{N}_{Z^{o}/Z}) , \]
qui est localement libre de rang n. Donc dans ce cas, les hypothèses préli-%
minaires faites plus haut sont vérifiées (j'ai oublié de dire qu'on supposait
X encore lis se sur Y), \over{D}{\unsure{d}}e plus on ne dispose pas\inserted{\unsure{,}} à priori d'un homomorphisme
surjectif $u:\underline{E}\rightarrow\underline{J}$ global, avec $\underline{E}$ loc libre de rant m sur X, mais en revanche
on dispose de l'homomorphisme global surjectif
\[ (\unsure{$\nmid$|$\dagger$}) \qquad \underline{E}^{o}\longrightarrow \underline{N}_{\xout{\illegible{1}}} \quad , \]
qui localement peut être considéré comme déduit d'un homomorphisme sur\over{g}{j}ectif
$\underline{E}\rightarrow\underline{J}$ par restriction à $V^{o}$ . Alors $V^{1}$, $V''$, \inserted{\unsure{$V^{2}$}}\unsure{p}euvent se définir directement
en termes de (\unsure{$\nmid$|$\dagger$}) , ainsi que $\underline{K}$ , $\underline{W}$ , de sorte qu'on trouve directement un
\[ \underline{M} : \underline{K}\rightarrow \underline{\mathrm{Sym}}^{2}(\underline{W}) \qquad\qquad \text{sur } V^{1}-V'' \]
intrinsèque à la situation, d'où \xout{\unsure{disc}} un discriminant
\[ D \in \Gamma(V^{1}-V'',\ \underline{K}^{-n+m-1}\hand{$\otimes$}\ \text{dét}(\underline{W})^{\hand{$\otimes$}2}\ ), \]
tel que $V(D) = V^{2}$ .

Je propose de nous contenter d'ajouter toutes ces considérations d'inva-%
riance en petits caractères, après le cas des sections de Modules loc libres;
je te laisse le choix de donner ou non une terminologie de transversalité dans
le dernier cas envisagé. Il reste la question dans quel paragraphe doit être
la place du fourbi supersingulier. Du point de vue résultats, le plus naturel
semblerait le par.19, mais du point de vue technique, dans les pages 18,19
j'ai utilisé un sorite différentiel qui, \xout{w}semble-t-il, viendrait plutôt au par
22 des grassma\unsure{nniennes}. Si tu ne remplaces cette démonstration savante par un
calcul bébête de matrices (certainement possible, mais moins instructif) il
semblerait qu'on devra rejeter la supersingularité au par.22.

\margin{171}
